\chapter{Συζήτηση}
\label{ch:discussion}

\guide{Έκταση 4--6 σελίδες.}

\section{Απαντήσεις στα ερευνητικά ερωτήματα}
\label{sec:answers}
\guide{Μία παράγραφος για καθένα από τα ΕΕ1--ΕΕ4 της \cref{sec:rq}, με τους αριθμούς του \cref{ch:evaluation}.}

\section{Περιορισμοί}
\label{sec:limitations}
\begin{itemize}
  \item Η ευθυγράμμιση εκτιμά στροφή μόνο γύρω από τον κατακόρυφο άξονα· προϋποθέτει σωστή κατακόρυφο από τη συσκευή.
  \item Η ακρίβεια περιορίζεται από το πόσο καλά ακουμπά ο χρήστης το σημείο με το χειριστήριο.
  \item Ο επανεντοπισμός των anchors εξαρτάται από τον φωτισμό και την υφή του χώρου.
  \item Ο ενσωματωμένος εξυπηρετητής προορίζεται μόνο για τοπικό δίκτυο (HTTP με κλειδί API).
  \item Η αναγνώριση φωνής εξαρτάται από υπηρεσία νέφους (Wit.ai).
  \item Η εφαρμογή χρησιμοποιεί το Built-in Render Pipeline της Unity και το SDK της Meta, άρα δεν μεταφέρεται αυτούσια σε άλλα headsets.
\end{itemize}

\section{Απειλές εγκυρότητας}
\label{sec:validity}
\guide{Μικρό δείγμα χρηστών, ένα εργοτάξιο, ένας τύπος συσκευής, ο ερευνητής είναι και ο δημιουργός του συστήματος (μεροληψία), εξοικείωση συμμετεχόντων με VR.}
